\note{sa page 13}

appliquant $L_{0}$ en $L$, et les $\pi_{i}$ en les $\pi_{i}$
\struck{\ill{}} \add{(des $\pi_{i} \in \mathcal{E}$)}. \struck{exprimant la graduation de
$\mathcal{E}$ par (1.3).} De plus, sous ces
conditions, l'hom.\ \struck{\ill{}} dans (ii) est unique, et avec les notations du n° 6,
il transforme $\Lambda_{0}$ en $\Lambda$, les $\pi_{i}$ en $\pi_{i}$, les
$\pi_{i,\alpha}$ en $\pi_{i,\alpha}$, les $\varphi^{\beta,\alpha}_{i}$ en
$\varphi^{\beta,\alpha}_{i}$ \struck{\ill{}}.

\marginal{condition essentielle !}

\emph{Dém.} (i) $\Rightarrow$ (ii) : Résulte de la table de
multiplication des $\varphi^{\beta,\alpha}_{i}$ donnée au n° 6, et du th.\
\ill{} donnée au n° 6.

\note{De « soient » à « Je dis que $\pi'_{i} = \pi_{i}$ … », le passage est encadré
et traversé de traits obliques : barré.}

(ii) $\Rightarrow$ (i) : soient \ill{} $\pi'_{i}$, \ill{}
$L^{\alpha} \pi'_{i} = \pi'_{i+\alpha} L^{\alpha}$ \ill{} de $\mathcal{E}$
\ill{}. Je dis que $\pi'_{i} = \pi_{i}$ pour tout $i$. Comme $i \leqslant n$, on a
\[ \bigl( \Lambda_{0}^{n-i} L_{0}^{n-i} - 1 \bigr)\, \pi_{i} = 0 , \qquad \bigl( L_{0}^{n-i} \Lambda_{0}^{n-i} - 1 \bigr)\, \pi_{2n-i} = 0 , \]
on en conclut, appliquant $\varphi$ (et pour $\Lambda' = \varphi(\Lambda)$)
\[ \bigl( \Lambda^{n-i} L^{n-i} - 1 \bigr)\, \pi_{i} = 0 , \qquad \bigl( L^{n-i} \Lambda^{n-i} - 1 \bigr)\, \pi_{2n-i} = 0 , \]
Donc la condition (ii) du th.\ du n° 6 est \uncertain{vérifiée}, donc aussi la
condition (i). \struck{\ill{}} \ill{} Pour l'unicité de $\varphi$, \ill{} de
voir que $\varphi(\Lambda) = \Lambda$, et
\struck{$\varphi(\pi_{i,\alpha}) = \pi_{i,\alpha}$} \add{simultanément} \struck{\ill{}}.
\ill{} la \uncertain{caractérisation} de $\Lambda$ et des $\pi_{i,\alpha}$ par les
formules (6.14), \ill{} $\varphi(\pi_{i}) = \pi_{i}$.

